problem:  
Let \((M^{2n+1}, T^{1,0}M, \theta)\) be a compact, strictly pseudoconvex CR manifold, and let  
\[
\mathcal{S}_f = \{u>0 \mid R_{u^{2/n}\theta} = f \}
\]
denote the set of conformal factors \(u\) solving the CR Nirenberg problem for a fixed smooth function \(f\).  

Assume \(f>0\) and normalized so that \(\int_M f\,\theta \wedge (d\theta)^n = \text{Vol}(M)\).  
It is known (Jerison–Lee, Gamara–Yacoub) that on \(S^{2n+1}\) with its standard CR structure, sequences of solutions may blow up along orbits of the CR automorphism group, and such bubbling limits correspond to Heisenberg solitons.  

**Problem:**  
Formulate and determine whether the following *compactness dichotomy* holds for general strictly pseudoconvex CR manifolds:

> (**Conjectural Compactness Principle**)  
> For \((M, T^{1,0}M, \theta)\) compact, the solution set \(\mathcal{S}_f\) of the CR Nirenberg problem is compact in \(C^{2,\alpha}(M)\) if and only if the CR automorphism group \(\mathrm{Aut}_{CR}(M)\) is compact.  
> 
> In the noncompact case, any sequence \(u_i \in \mathcal{S}_f\) that fails to have a convergent subsequence must, after applying elements \(\phi_i \in \mathrm{Aut}_{CR}(M)\) and rescaling, subconverge to a nontrivial entire solution of the CR Yamabe equation on the Heisenberg group.  
> 
> **Question:** Is the above dichotomy universally valid for all compact strictly pseudoconvex CR manifolds admitting a pseudo-Einstein contact form?

Why is it a "good" problem:  
This question explicitly proposes a **CR compactness dichotomy conjecture**, paralleling a central structure result in conformal geometry but unproven in the CR setting. It removes trivialities (such as automatically bounded curvature) and directly asks whether the *geometry of the automorphism group* is the unique obstruction to compactness of the solution space—an open and fundamental issue. It bridges **geometric analysis (nonlinear PDE compactness and blow-up analysis)** and **complex geometric structures (CR automorphism symmetry)**, seeking a universal principle relating symmetry to analytic compactness. Its simple formulation belies deep implications: confirming or refuting it would unify known phenomena (spherical bubbling, Heisenberg limits) into a global structural theorem, potentially foundational for the CR analogue of the global Nirenberg theory.